\chapter[Love, Utility, Law, and War]{Love, Utility, Law, and War: How the Hundred Schools Argued}
\section{A Worn-Out Pair of Hemp Shoes}
Begin with an object: a pair of hemp shoes.

Not trainers in someone's cupboard, but the kind lying behind museum glass: woven from hemp cord, soles worn through, half the uppers unravelled, greyish yellow after more than two thousand years. Similar shoes have emerged from Warring States burials. Most of their owners left no name behind.

Today, however, imagine that they belong to someone whose name we know. In the late fifth century BCE, a man called Mo Di wore such shoes as he left Lu and walked for ten days and nights. When the soles wore through, he tore strips from his clothes, wrapped his feet, and kept walking until he reached Ying, the capital of Chu. He was not fleeing disaster, conducting business, or seeking relatives. He had come to interfere in something that was none of his concern: Chu intended to attack Song, and he intended to stop it.

Song was not his homeland. None of the people there waiting to die knew him. His journey promised neither payment nor office, and bandits or disease might kill him at any point.

Do not rush to feel moved; first, find it strange. Why would someone wear out a pair of shoes for strangers? Emotion comes cheaply; puzzlement is valuable. Behind this strangeness stands an entire view of the world, and Mo Di was not alone in having one. In China at that time, some proposed saving the world through love; others argued that love caused the trouble and only impersonal institutions were reliable. Some probed language, disputing whether a white horse counted as a horse. Others studied warfare, only to discover that the crucial question was not how to fight but how much you knew before fighting.

This chapter enters that great dispute, later called the ``contention of the Hundred Schools.'' Remember the worn-out shoes: they carried one of its arguments farther than any others.

\section{A Military Exercise below Ying's Walls}
Come into the scene.

The time: a hot, dry afternoon around 440 BCE. The place: near the palace at Ying, Chu's capital, around present-day Jiangling in Hubei. Dust and bronze hang in the air. Craftsmen have been rushing to build a new siege device: a scaling ladder that can be pushed to a wall for soldiers to climb. Its inventor is Gongshu Ban, later honoured as Lu Ban, patron of carpenters, and the greatest engineer of his day.

Two men stand on open ground below a gate tower. One is Gongshu Ban, carrying his ladder plans inside his robe. The other is Mo Di, who has walked for ten days and nights: ragged clothes, torn fabric wrapped round his feet, blood mixed with dust.

``Gongshu'' in the \textit{Mozi} records what followed. Mo Di first visited Gongshu Ban, then the king of Chu, but reasoning failed. The king said: the ladders are already built; why not attack? Mo Di proposed a trial. He removed his belt, curved it into a circle on the ground to represent Song's walls, and used small pieces of wood as attacking and defensive devices. Gongshu Ban attacked; Mo Di defended.

This may have been one of military history's earliest war games. Gongshu Ban changed his method nine times; Mo Di repelled him nine times. Gongshu Ban exhausted his siege devices, while Mo Di still had defensive methods to spare.

Having lost the exercise, Gongshu Ban suddenly said something chilling: ``I know how to deal with you, but I will not say.''

Mo Di replied: ``I also know how you intend to deal with me, and I will not say either.''

The king asked what they meant. Mo Di explained: Gongshu Ban simply wants to kill me. But three hundred disciples, including Qin Guli, are already on Song's walls with my defensive equipment. Killing me will not kill them all.

The king was silent for a while, then said: Very well; I shall not attack Song.

Before it began, a war was cancelled by hemp shoes, a belt, and three hundred students stationed on city walls.

Notice the three things that mattered. First, Mo Di's eloquence---which failed to persuade the king. Second, the exercise---which showed that attacking would cost more than it was worth. Third, the three hundred disciples---who proved that behind Mo Di stood a disciplined group prepared to die. Reasoning, calculation, organisation: only with all three did the war stop. The rest of this chapter examines how the Hundred Schools separately honed these into weapons.

\section[Where Does Order Come From?]{The Real Question in an Age of Disorder: Where Does Order Come From?}
To understand the dispute, we must know what was at stake. Lay out the conflict before rushing to an answer.

In Mo Di's lifetime, the old world was collapsing. The realm enfeoffed by the Zhou king had rested on two things: kinship, as he granted land to relatives and meritorious followers, and ritual propriety, rules specifying who should do what and how. By the Spring and Autumn and Warring States periods, neither worked. Three ministerial houses, Han, Zhao, and Wei, partitioned Jin; the Tian clan replaced Qi's Jiang rulers. Ministers dividing their lord's state was monstrous disloyalty under the old rites, yet it happened. Warfare no longer observed ``honourable battle formations''; mass killing of prisoners, sieges, and the destruction of states became commonplace. The floor of the old order collapsed. Everyone was falling, and everyone asked the same question:

\textbf{On what can a new order be built?}

Underneath lay a sharper question: can human nature be relied upon?

If it can---if people already possess love and an inability to bear others' suffering---salvation means cultivating and enlarging these. If it cannot---if people naturally seek benefits and avoid harm, responding to rewards and behaving only under punishment---moral preaching wastes time. Better design a machine of clear rewards and penalties. The question survives today. Should school discipline depend on ``collective pride'' or a deductions chart? Should companies rely on organisational culture or performance assessments? Two ghosts from over two thousand years ago stand behind each dispute.

Warring States answers roughly followed four routes, the four schools examined here. These are not four slogans but four complete prescriptions, each containing a diagnosis, treatment, and side effects:

\begin{itemize}
\item \textbf{Mohists} said: the disease is caring only for one's own people. The remedy is ``inclusive care,'' extending love to strangers, together with ``honouring the worthy'' (promoting capable people), ``frugality'' (avoiding waste), and ``opposing aggression.'' Mo Di's hemp shoes put this prescription into motion.
\item \textbf{Legalists} said: the disease is expecting morality from people at all. The remedy is ``law'': explicit rewards and punishments, combined with the ruler's power and techniques. Whatever people's hearts are like, prevent bad people from carrying out bad deeds.
\item \textbf{The School of Names} said: stop arguing for a moment. What exactly do ``love,'' ``law,'' and ``righteousness'' mean? If you cannot even explain whether a white horse is a horse, how dare you discuss government? First calibrate language, your measuring instrument.
\item \textbf{Military thinkers} said: while you argue, the enemy is moving troops. First acknowledge the world's dangers, then learn to decide amid incomplete information and changing circumstances.
\end{itemize}
The four routes scorned one another. Mohists condemned war; military thinkers studied it. Legalists mocked morality; Mohists built their lives upon it. All three mocked the School of Names for ``word games,'' yet none could evade its questions. This was no agreeable tea party but a real fight, with each side holding questions the others could not answer.

Before continuing, cast a private vote: which prescription would you stake your hopes on to rebuild a shattered order?

\section[Two Calculations of Love]{Two Calculations of Love: Inclusive Care against Graded Affection}
Hear the Mohists out first, then their opponents.

Mohism is a complete package; we cannot select only its pleasant-sounding pieces. It begins with ``inclusive care.'' The first ``Inclusive Care'' chapter of the \textit{Mozi} ends with:

\begin{quote}
Thus when everyone under heaven cares inclusively for one another, there is order; when they hate one another, there is disorder.
\end{quote}
In essence: mutual love brings stability; mutual hatred brings chaos. Mo Di diagnoses disorder as \textit{bie}, partiality: people love only their own and treat others' households and states as expendable. Ministers disrupt other households to enrich their own; rulers attack other states to enlarge theirs. The logic is a thief's, merely on another scale. Replace partiality with inclusiveness: love others' parents as your own and others' states as your own.

You probably want to object: how could anyone do that? Wait; the Mohists are ready. They attach an engine called ``mutual benefit.'' Loving others is not unrewarded: care for them, they care for you, and everyone gains; harm them, they harm you, and everyone loses. Inclusive care is not merely a noble feeling but an account that balances. This is Mohist ``utility'': benefit for everyone under heaven. The Chinese term often sounds disparaging today, suggesting calculating self-interest; for Mohists it is profoundly serious. Whether something is right depends on the benefit or harm it brings the world.

This engine drives the other doctrines. ``Honouring the worthy'': if the world is to be governed well, the most capable should rule, not the ruler's brother-in-law. The first ``Honouring the Worthy'' chapter states boldly:

\begin{quote}
Officials need not remain forever noble, nor common people forever lowly.
\end{quote}
In other words, officeholders will not always be exalted, nor ordinary people always humble. Over two millennia ago, when hereditary aristocracy seemed self-evident, this nearly amounted to rebellion. ``Frugality'' and ``moderation in funerals'': Confucians advocated three years' mourning and lavish burial; Mohists counted the cost. Production stops, bodies weaken, treasures go underground. Who pays? Everyone. Then ``opposing aggression'': opposition to wars attacking other states, not to all warfare. Mohists themselves were the world's best-known defensive experts. Disciples trained hard in city defence and went wherever aggression threatened. It is easy to misread this: treating them as passive pacifists who never resist is mistaken. They distinguished ``aggression'' from ``punitive action.'' Invasion was unjust; resistance and campaigns against tyrants were different matters. The shoes walked towards Ying not to persuade Song's people to submit to slaughter, but to make invasion unwinnable.

One often overlooked fact remains: Mohism was an organisation as well as a school. It had a leader, the \textit{juzi}, and iron discipline. The late Warring States \textit{Spring and Autumn Annals of Master Lu} tells of a leader called Fu Tun whose son killed someone. Because Fu Tun was old and had only one son, Qin's king pardoned the son from execution. Fu Tun replied that Mohist rules required ``death for killing, punishment for injury,'' and he could not fail to enforce them. He ultimately executed his own son under the group's discipline. This story comes from another school's writings, and not every detail is necessarily trustworthy. But its image was widely accepted: Mohists placed discipline above family affection.

Now hear the opponents, whose reasoning we shall state fully, even more clearly than they did themselves.

Confucianism was Mohism's lifelong adversary. Mencius attacked Mo Di fiercely: inclusive care denied a father's special status and was almost beastly. Cool that heated language into a discussable argument. Confucians advocated ``graded affection'': love begins with those nearest us and extends outwards in circles. Loving parents more deeply than passers-by is not selfishness but love's normal shape. Their full case has at least three layers:

\begin{itemize}
\item First, love is an ability, not a switch. Kindness to strangers develops through loving particular people nearby. Someone cold towards the parents they live with yet claiming equal love for humanity probably loves the uplifting concept ``humanity,'' not any actual person.
\item Second, gradation grounds responsibility. Because parents have given me special care, I owe them a special debt. Flatten love and you flatten responsibility; flattened responsibility becomes no responsibility. Telling a stranger ``I love you equally'' cannot cancel what you owe your family.
\item Third, impossible demands mass-produce hypocrisy. Set the standard at ``love passers-by as your mother'' and nobody can meet it, so everyone performs. Should moral standards be as high as possible or set where they can be fulfilled? This is a real question.
\end{itemize}
Mohists can answer: those circles license war and favouritism. ``My country first'' is ``my family first'' enlarged. Confucians can reply: love without circles is either impossible or imposed through discipline. Look at the Mohists: do they rely on spontaneous inclusive care or the leader's orders? What happened to love?

\section[Institutions and the Fog of War]{Two Cold Calculations: Legalist Institutions and the Military Fog}
Mohists wager that people can be persuaded to love. The next two schools wager differently: trust calculation, not the human heart.

First, Legalism. Its great synthesiser was Han Fei, a late Warring States prince of Han. He stammered and lacked eloquence, so put his words into writing. Before him, Legalism possessed three foundations: Shang Yang's ``law,'' public provisions with clear rewards and punishments; Shen Buhai's ``technique,'' the ruler's privately held methods of controlling ministers; and Shen Dao's ``positional power,'' authority conferred by office---people obey your position, not your virtue. Han Fei assembled them into one machine: law supplies public tracks, technique the hidden steering, and positional power the engine. Without any one, the vehicle cannot move.

Legalism's assumptions about human nature are famously cold. ``Precautions within the Palace'' in the \textit{Han Feizi} says in effect: carriage makers hope everyone becomes wealthy; coffin makers hope everyone dies young. Not because coffin makers are wicked, but because if nobody dies, coffins do not sell. The original states:

\begin{quote}
When a craftsman makes coffins, he wishes people to die young.
\end{quote}
Notice the point. Han Fei is not condemning human wickedness. He says behaviour follows interests as water flows downhill, independently of goodness or badness. Government should therefore design reward and punishment---his ``two handles''---rather than educate people into goodness. Let self-interested people do what is required by following their interests. Moralists want to change the world; Legalists want to calculate it.

This approach has two impressive achievements. First, it opposes privilege. ``Having Standards'' in the \textit{Han Feizi} says:

\begin{quote}
Law does not flatter the noble; the marking line does not bend for crooked wood.
\end{quote}
In essence: law does not favour the powerful, just as a carpenter's ink line does not accommodate bent timber. Nobles and commoners receive the same punishment for offences---a startling claim in an aristocratic age. Second, it opens paths for ordinary people: military merit earns rank, good farming earns exemption from labour service, and birth matters less. Qin's transformation from a western border state into a war machine through these institutions was no accident.

But stop at a common misjudgement. Reading ``law does not flatter the noble,'' many immediately think of today's ``equality before the law'' and imagine Legalism pioneered the rule of law in ancient China. Wrong. Modern rule of law places law above everyone, including the supreme ruler. Legalist law is the ruler's tool; he stands above and outside it. Law is a whip, and the hand holding it is not whipped. One places the king under law; the other governs by wielding law as a whip. The directions are opposite. A small verbal distinction spans two millennia. This is not pedantry: countless online disputes collapse by confusing the two.

To give Legalism its full case, add its strongest challenge. Han Fei asks: Confucians and Mohists have preached morality for two centuries, yet wars grow larger and deaths multiply. Where have your proposals delivered? ``The Five Vermin'' summarises the age:

\begin{quote}
In high antiquity people competed in virtue; in the middle age they pursued stratagems; today they contend in force.
\end{quote}
In essence, ancient people competed morally, those of the middle age through clever plans, and contemporaries through strength. The same chapter tells a famous fable. A farmer in Song has a stump in his field. A running hare strikes it and dies. Delighted by the free meat, he abandons his tools and waits beside the stump. No more hares arrive, and his field goes to waste. Governing today's world with ancient sage-kings' methods, Han Fei says, is waiting beside a stump for another hare. This is Chinese political thought's most famous hare.

Legalist calculation faces inward. Another calculation faces outward: military thought.

Its founding figure, Sun Wu, supposedly presented thirteen chapters of strategy to Wu's king in the late Spring and Autumn period. The most counterintuitive feature of \textit{The Art of War} is that a book about war opens not with fighting but with avoiding miscalculation. Its central action precedes battle: compare politics, weather and timing, terrain, commanders, and regulations; assess the chances of victory; do not fight if the calculation does not favour you. ``Disposition'' states:

\begin{quote}
Thus victorious armies first secure victory, then seek battle; defeated armies first fight, then seek victory.
\end{quote}
In essence: capable armies first create a winning position, then fight; incapable ones rush in and hope to win by luck. War is not an arena of courage but an information problem. The best-known line in ``Planning Attacks'' says:

\begin{quote}
Know the enemy and know yourself, and a hundred battles will bring no peril.
\end{quote}
Information is always incomplete, so military thinkers developed methods for handling uncertainty: creating false appearances. ``Calculations'' says ``Warfare is the way of deception'': when capable, appear incapable; when acting nearby, appear to be acting far away. The final chapter, ``Using Spies,'' deals entirely with espionage. The battlefield is fog; whoever clears some of the enemy's fog while throwing more sand in the enemy's eyes is halfway to victory.

Put Legalists and military thinkers together and their shared worldview emerges: separate ``what ought to be'' from ``what actually is,'' inspect the cards before discussing ideals. This is not uniquely Chinese. In his \textit{History of the Peloponnesian War}, Thucydides records Athenian envoys telling the besieged Melians, in essence, that the strong do what they can and the weak endure what they must. India's \textit{Arthashastra} describes interstate relations through the ``law of the fish'': large fish eat small ones, and a ruler's first lesson is not to be a small fish. The Warring States period was no exception, but a recurring human predicament: when rules fail, calculation takes over.

\section[Thinking Bridge: Distinguish Concepts]{Thinking Bridge: Names and Realities---Taking an Argument Apart through Conceptual Distinctions}
After all this argument, step back and take away a tool: how to dismantle any dispute for inspection. It is called conceptual distinction, and its Warring States supplier was the School of Names, mocked by every other school.

Its signature problem was Gongsun Long's ``A white horse is not a horse.'' His ``Discourse on the White Horse'' opens as a dialogue:

\begin{quote}
``A white horse is not a horse''---can this be affirmed? The reply: It can.
\end{quote}
How can that hold? His argument breaks down like this. Ask for a horse, and a yellow or black horse will do. Ask for a white horse, and neither will do. Since ``horse'' and ``white horse'' respond differently to the same request, they are not the same concept. Moreover, ``horse'' names a shape, ``white'' a colour; ``white horse'' combines colour and shape, which plainly differs from shape alone.

You may instinctively laugh: surely this is mere contrariness? A white horse is obviously a horse and sells as one at the horse market. Hold back: this is the chapter's first ``easily misjudged counterexample.'' Collections of historical absurdities treat it as the ultimate pedantry, but it asks a real question: \textbf{by what standard do we determine relations of inclusion between concepts?} ``A white horse is a horse'' concerns classification: white horses belong to the class of horses. Gongsun Long's ``A white horse is not a horse'' concerns identity: the two concepts are not identical. Both statements can hold; the quarrel arises because the same ``is'' performs two jobs. Belonging is not equality. Gongsun Long lacked modern logical symbols, but found that gap.

Another branch, represented by Hui Shi, moved oppositely: Gongsun Long divided energetically; Hui Shi united energetically. He advocated ``uniting sameness and difference'': both are relative. From one perspective, horses and cattle are four-legged animals; from another, even the same horse differs between yesterday and today. One cuts the world apart, the other blends it together. Combined, they offer a complete language check-up: where does each word's boundary lie, and does it survive a change of setting?

Greece had a similar worry at roughly the same time. Travelling teachers called Sophists taught courtroom argument, specialising in turning an opponent's ambiguous words against them. Protagoras's famous saying means, broadly, ``Man is the measure of all things.'' Greeks soon saw the danger: push conceptual games far enough and the cleverest tongue wins while truth is forgotten. Anxiety that language might hold thought hostage was an epidemic of the age across humanity.

Now assemble the tool. Analyse any dispute in three steps:

\begin{itemize}
\item \textbf{First, definitions}: what does each side mean by the keywords? (``Rule of law'' in Han Fei and in a modern constitutional textbook refers to different things.)
\item \textbf{Second, boundaries}: what does the concept include or exclude? (Does ``love'' include enemies? Does ``war'' include cyberattacks?)
\item \textbf{Third, counterexamples}: find a borderline case and force each side to respond. (White horses, hares, tomatoes as fruits or vegetables.)
\end{itemize}
Add a companion move: draw an argument map. Divide any position into ``claim $\rightarrow$ reason $\rightarrow$ hidden assumption.'' For Mohism, the claim is ``We should practise inclusive care''; the reason, ``Disorder arises from the absence of mutual love''; the hidden assumption, ``People can love universally.'' Notice how two millennia of Confucian fire targets precisely that assumption. A dispute's real battlefield often lies not in the visible claim but in an unspoken assumption carrying everything. Next time you watch an online argument, ask what each side takes for granted without saying.

\section{An Anti-War Text from Two Millennia Ago}
Now read a short primary passage. After so much paraphrase, face the words of Mo Di, or his disciples, directly in the first ``Opposing Aggression'' chapter of the \textit{Mozi}. Your new tools fit its method perfectly.

It starts very small. Someone enters another person's orchard to steal peaches and plums. Everyone calls it wrong; officials punish the thief if caught. Why? Because the thief harms others for personal gain. The examples then escalate: stealing chickens, dogs, and pigs is worse; entering a livestock enclosure and taking horses or cattle is worse still; killing an innocent person and taking their clothes and weapons is worse again. Then comes the crucial statement:

\begin{quote}
Killing one person is called unrighteous, and necessarily constitutes one capital offence.
\end{quote}
In essence: killing one person is unjust and merits death. Follow the logic: killing ten is ten times as unjust, ten capital offences; killing a hundred is a hundred times as unjust, a hundred capital offences. Every gentleman reading along nods: yes, these are injustices deserving punishment.

Then the text reveals its blade:

\begin{quote}
Yet when it comes to the greatest unrighteousness, attacking a state, they do not know to condemn it; instead they praise it and call it righteous.
\end{quote}
In essence: attacking another country is the greatest injustice, yet people fail to oppose it and praise it as ``righteous.''

The text adds a superb analogy. Someone sees a little black and calls it black, but sees much black and calls it white. Everyone would say this person cannot distinguish the two. The original reads:

\begin{quote}
If one sees a little black and calls it black, but sees much black and calls it white, we would judge that person unable to distinguish white from black.
\end{quote}
Everyone recognises killing one person as unjust, yet calls killing ten thousand righteous. Is that not calling a little black black, but much black white?

Use the preceding section's tools. Claim: attacking a state is the greatest injustice. Reason: it shares the structure of theft and murder, merely on a larger scale. Hidden assumption: \textbf{moral judgements should not reverse when scale increases}. Stealing a peach makes a thief; stealing a country should not make a hero. Over two millennia later, that assumption remains among international news's hottest questions. Notice also that the text invokes no sage's quotation and makes no emotional appeal. It builds step upon step until you have nowhere to retreat. This is the Mohist style: the craftsperson's hands before the philosopher's mouth.

The short essay also offers a reading method worth copying. It never proves ``War is wrong.'' It proves: \textbf{if you think theft and murder are wrong, you cannot call attacking a state right}. It does not force a conclusion upon you; it checks your consistency. This is argument's hardest move to resist: I need not make you accept my premises; I defeat you with your own.

\section[Why the Schools Contended]{Why Did the Hundred Schools Contend? Three Explanations}
Pull back. One fact---the sudden emergence of dozens of conflicting doctrines during centuries of Spring and Autumn and Warring States history---admits genuinely different explanations. As usual, distinguish four things: what happened (schools proliferated, a matter of historical evidence); why (the competing explanations below); how contemporaries understood it (saving the world, not ``doing scholarship''); and how we evaluate it (calling it a ``cultural golden age'' is a later judgement, not a fact). We compare the second level.

\textbf{Explanation One: Society's reshuffling reshuffled thought.}

The case: hereditary aristocracy broke down, and commoners could enter power through military merit, eloquence, and ability. The legend of Su Qin holding the ministerial seals of six states captures such mobility. When old identities shattered, old principles required retelling. Knowledge had previously belonged to government, locked in official hands. Confucius opened private teaching without class distinctions, and learning flowed to ordinary people for the first time. More educated people meant more voices. This explanation connects intellectual history to social history and explains why contention arose then rather than in peaceful Western Zhou. Its limit: it explains the conditions for debate, not why these particular doctrines appeared. The same social change might produce only one school.

\textbf{Explanation Two: Warfare put ideas out to tender.}

The case: the seven great states fought year after year, and every ruler sought an answer to the same question: how to enrich the state and strengthen its army? Doctrines were bids. Mohists offered ``inclusive care, opposition to aggression, and city defence''; Legalists, ``farming, warfare, rewards, and punishments''; military thinkers, ``calculate victory before fighting.'' Qin adopted the Legalist proposal and unified the realm: the tender's result was announced. This sharp explanation accounts for the eventual success of Legalist and military thought and Mohism's rapid decline after unification. A unified empire no longer needed roving defensive organisations and instead saw them as threats. Its limit: reducing every idea to a useful bid cannot explain ``useless'' studies like the School of Names or internal disputes about details unattractive to rulers.

\textbf{Explanation Three: Division itself made a marketplace for ideas.}

The case: multiple states meant neither one authorised answer nor uniform censorship. A scholar unwelcome in Lu could pack up and move to Qi. Qi's Jixia Academy supported scholars from many states with mutually hostile positions, allowing them to ``discuss affairs of state without holding office'': argue without taking responsibility. Like products, ideas developed through competition among states. This explains a counterfactual question: why did contention end after unification? Qin's book burning and Han's ``exclusive honouring of Confucian learning'' became possible because the realm already acknowledged one supreme authority. Its limit: potential circularity---division explains contention, while contention proves division's vitality. Moreover, division and warfare are two sides of one situation, making this hard to separate from Explanation Two.

A broader account offers another reference point. The twentieth-century German scholar Jaspers proposed an ``Axial Age'': approximately 800--200 BCE, when China, India, Greece, and Israel independently produced foundational thinkers. It reminds us that the Hundred Schools did not sing alone in a sealed world. Why this happened simultaneously still has no agreed answer. Keep it as an open question, not a fact to memorise.

Each explanation holds part of the truth. An honest synthesis: social change supplied the people, war the questions, and division the meeting place. Without any one, the great debate could not have convened.

\section[Further Challenge: Weighing Benefit]{Further Challenge: Can We Weigh the Benefit of the Majority?}
Now for the chapter's weightiest theoretical problem, running from Mohist ``mutual benefit'' to today.

As discussed, Mohists judge right and wrong by ``benefit'': what benefits everyone under heaven is right. More than two millennia later, this found a famous echo across the Atlantic. Around the turn of the nineteenth century, the British thinker Bentham proposed judging actions and institutions by whether they produce ``the greatest happiness of the greatest number.'' This became known as utilitarianism. Its Chinese name can misleadingly suggest opportunism; it actually concerns ``utility,'' the total account of pleasure and pain.

Mohists and Bentham reach across two thousand years to shake hands: judge consequences, not noble motives or the antiquity of rules. This approach has enormous force, cutting through countless excuses that ``ancestral rules cannot change.'' Today's railway construction, health-insurance policy, and vaccine allocation all carry traces of this calculator.

Now oppose it and find where the calculator crashes. The following classic counterexample challenges the easy assumption that utilitarianism is unbeatable.

Suppose an innocent person lives in a city where everyone hates him. Publicly executing him would improve the mood of several million inhabitants. In a ``total happiness'' account, one sacrificed, millions pleased: a net gain. Should it be done?

Almost everyone's intuition shouts no. But notice: your objection now rests not on utility but on another principle---``The innocent must not be sacrificed''---which refuses the accounting. The counterexample shows that \textbf{aggregate calculations can swallow minorities}. The majority's happiness cannot automatically purchase permission to dispose of any particular person.

Return to Mohism. Its ``benefit'' takes the whole world as its unit and naturally favours the majority, yet ``inclusive care'' loves every individual. These pillars usually support one another but may conflict in extreme cases. If benefiting the world requires sacrificing a minority, what does inclusive care say? History preserves no answer from Mo Di, but preserves the Mohists' example. They spoke of ``entering fire and treading blades''; ``All under Heaven'' in the \textit{Zhuangzi} describes them as ``taking self-hardship to the utmost.'' Notice the direction: they asked \textbf{their own members} to sacrifice, not other people to balance the books. Is that an answer? Half of one. Voluntary and forced sacrifice differ; how voluntary ``voluntary'' remains in an iron-disciplined group is another question. That half is yours.

Utilitarianism has a deeper technical difficulty: \textbf{what is the unit of happiness?} How do you convert your pleasure in ten extra exam marks into a friend's pleasure at avoiding tutoring? Mozi was at least concrete about ``benefit'': enough food, no war, fewer deaths. But concrete accounts provoke further questions: what after everyone has enough to eat? Convert all life into ``benefit'' and the person disappears into the conversion.

This book answers none of these questions. But keep the calculator: whenever someone says ``It is for everyone's good,'' ask who ``everyone'' includes and what happens to the person left out.

\section{Today: The Philosophers in Class}
This twenty-three-century-old quarrel is holding a local session beside you right now.

\textbf{Debates over class rules replay Confucian--Legalist debates.} A teacher has two disciplinary routes: shared identity and voluntary responsibility---``Our class is one big family''---or conduct marks: two deducted for lateness, three for skipping cleaning duty, end-of-term honours based on the total. You have probably seen both fail. Rely on feeling, and someone treats collective goodwill as a free lunch. Rely on deductions, and someone soon makes ``not getting caught'' their chief skill. Han Fei would say: see, do not rely on conscience. Mohists would ask: is a classroom reduced to a deductions chart somewhere you want to be? Your class resits this test monthly; what is missing is your answer.

\textbf{Examination selection inherits both ``honouring the worthy'' and its difficulty.} ``Officials need not remain forever noble, nor common people forever lowly'' eventually grew into imperial examinations, then today's gaokao: selection by marks rather than fathers. Mohists dreamed of it, Legalists partly realised it, and more than a millennium refined this enormous mechanism. Living today, Mohists might ask who defines worthiness. Marks measure endurance at practice questions, but do they measure the ability to improve a community? When honouring the worthy offers only one track and everyone crowds onto it, we get the familiar \textit{neijuan}, exhausting competition. Mohists invented the problem, not its escape route.

\textbf{``A white horse is not a horse'' appears daily in your group chats.} In a forwarded screenshot, one person says ``I was joking,'' another ``That is bullying''; one says ``self-defence,'' another ``a mutual fight.'' Under every dispute lies a conceptual boundary. Where does joking become bullying? At what step does defence become fighting? The School of Names is immediately useful: before choosing sides, ask whether both parties mean the same thing by their keywords. Half of online abuse begins when one word means two things; the other half when someone deliberately exploits that fact.

\textbf{Military thinkers inhabit every fog of information.} Before scheduling revision, you study past papers and your weaknesses: ``secure victory before seeking battle.'' In a game, experts inspect the minimap and estimate the enemy jungler's location before hiding in the grass: ``know the enemy and yourself.'' An influencer's carefully engineered ``casual clip'' exemplifies ``when capable, appear incapable.'' Sunzi today might work in risk management or intelligence analysis. His most important lesson is also unfashionable: do not fight battles the calculation says you cannot win. But engagement rewards ``charge first, think later.'' ``Fight first, then seek victory'' has become the internet's default posture.

\textbf{Reread ``Opposing Aggression'' alongside the news.} Whenever a large state attacks a small one, rhetoric casting aggression as justice arrives on schedule: protecting nationals, pre-emptive action, special operations. Mo Di's stepwise argument has lost none of its force. Between individuals, forcibly entering a home is criminal; between states, how does it become righteous? ``Call a little black black, but much black white'': that ancient mockery sounds like yesterday's trending news.

\textbf{The ``everyone's good'' calculator is now mass-produced.} ``Give up this opportunity for the class's honour.'' ``Can your problem wait so everyone else's revision is undisturbed?'' Each time, apply this chapter's questions: who is ``everyone'' in the account, and who lies outside it? Is the concession voluntary, or squeezed out by the word ``collective''? Mohists spent their lives showing that sacrificing for others can be noble. They authorised nobody to announce someone else's sacrifice.

The Hundred Schools never left. They merely changed classrooms.

\section[Your Turn: Play One of Three Cards]{Your Turn: Three Cards, but You May Play Only One}
Return to the worn-out shoes. Put the ancient problem into a situation you might actually face.

A new student joins your class. Half a term later, he is the excluded one: nobody wants him in a group, invites him into chats, or sits with him at break. No punches, no insults---the isolation stays below the surface. When the teacher asks, everyone is innocent: ``We haven't bullied him.''

The teacher decides to address it seriously and asks you to advise. Three cards lie on the table. You may play only one:

\begin{itemize}
\item \textbf{The Mohist card}: launch an ``inclusive care'' campaign. Encourage everyone to approach him, include him in cleaning duties, groups, and activities, and dissolve isolation through goodwill. Cost: campaigns often blow over quickly, and kindness that others are ``required'' to offer may feel uncomfortable to receive.
\item \textbf{The Legalist card}: establish a written class rule. Proven exclusion or isolation of a classmate incurs deductions, including for the offender's group. Cost: isolation mostly happens beyond the reach of evidence. A rule that controls actions but not looks may drive exclusion further underground.
\item \textbf{The School of Names card}: do not act yet. First have the class define ``bullying.'' What counts as exclusion, and what as the freedom not to play together? Could accusations without clear boundaries wrong innocent people? Cost: definition debates may become a talking shop while the student still sits alone.
\end{itemize}
There is no fourth card. The rule is one card only. Choose and give your reasons.

Before writing, remember what you already hold. Mohists tell you that a bystander's ``doing nothing'' is itself a kind of action. Legalists tell you that toothless goodwill will be consumed by whoever cares least. The School of Names tells you to calibrate words before accusing, lest justice strike the wrong person. Military thinkers add: find out what this student wants before playing your card. Even doing good cannot escape ``know the other and yourself.''

Remember, too, that there is no standard answer. Supporters of all three cards can find ancestors among the Hundred Schools, and two millennia of argument produced no unique solution. Your only difference from them is that this time, the person sitting on Song's walls waiting for the outcome may be you.
